\section{Natural Numbers}
\label{sec:naturali}

Although the properties of natural numbers should be intuitively well-known,
their formalization is not trivial. In particular, some proofs will be
complicated because when we need to prove many obvious facts, we must be careful
not to use properties that seem obvious to us but have not yet been proven. For
this reason, it is perfectly legitimate to skip the proofs on a first reading to
understand the general sense of the discussion and then, if necessary, go back
to delve into the details.

The natural numbers $0,1,2,\dots$ are the numbers we use for counting or for
iterations. There is a first natural number, which for us will be $0$, and then
for every natural number $n$, there is a successor which we will call
$\sigma(n)$. Starting from $0$ and moving to the successor, all natural numbers
are reached.

\begin{definition}[Peano's axioms]
  \label{def:assiomi_peano}%
  \label{def:naturali}%
  \mynote{Giuseppe Peano (1858--1932) see historical notes on page~\pageref{nota:Peano}} 
  \index{numbers!natural}%
  \index{$\NN$}%
  \index{Peano}%
  \index{axioms!of Peano}%
  \index{set!inductive}%
  \index{inductive}%
Given a set $\NN$, 
an element $0\in \NN$, 
and a function $\sigma\colon \NN\to\NN$, 
we will say that $(\NN,0,\sigma)$
satisfies Peano's axioms,
or more simply that $\NN$ is a set of natural numbers,
if the following properties hold: 
\mynote{The function $\sigma$ gives us the successor of each number. 
We will define $1=\sigma(0)$, $2=\sigma(1)$, $3=\sigma(2)$, \dots
The first Peano axiom states that different numbers have 
different successors. 
The second axiom tells us that the first 
natural number, for us $0$, 
is not the successor of any other number while every other number 
has a predecessor. The third axiom, the principle of induction, 
tells us that every natural number can be reached from $0$ 
by making a finite number of steps to the successor.}
\begin{enumerate}
  \item $\sigma$ is injective;
  \item $\sigma(\NN) = \NN \setminus\ENCLOSE{0}$;
  \item if $A\subset \NN$ and 
  \begin{enumerate} 
    \item[(i)] $0\in A$ 
    \item[(ii)] $n\in A \implies \sigma(n)\in A$
  \end{enumerate}
  then $A=\NN$.
\end{enumerate}
\end{definition}

From now on, when we use the symbol $\NN$, we will always assume that $\NN$ is a
set that satisfies the previous definition. We will see in
Theorem~\ref{th:unicitaN} that two different sets that satisfy Peano's axioms
are isomorphic, meaning they can be put in a one-to-one correspondence while
maintaining the same structure. In this sense, we can say that Peano's axioms
uniquely characterize the set of natural numbers. Therefore, we will say that
$\NN$ is \emph{the} set of natural numbers rather than \emph{a} set of natural
numbers.

Given $n\in \NN$, the number $\sigma(n)$ is called the \emph{successor}%
\mymargin{successor}%
\index{successor} of $n$. 
The number $0\in \NN$, which by axiom is not the successor of any other 
natural number, is called \emph{zero}%
\mymargin{zero}%
\index{zero}. 
For convenience, we also give names to the other \emph{decimal digits}
\mymargin{decimal digits}%
\index{digits!decimal}%
that is, the first natural numbers
(later we will introduce positional notation to represent every natural number 
with a sequence of decimal digits):
\begin{equation}\label{eq:cifre}
\begin{gathered}
 1 \defeq \sigma(0),\quad  
 2 \defeq \sigma(1),\quad
 3 \defeq \sigma(2),\quad 
 4 \defeq \sigma(3),\quad
 5 \defeq \sigma(4),\\ 
 6 \defeq \sigma(5),\quad 
 7 \defeq \sigma(6),\quad 
 8 \defeq \sigma(7),\quad 
 9 \defeq \sigma(8)
\end{gathered}
\end{equation}
and we observe that $\sigma$ is the usual counting operation:
 \[
 0 \stackrel\sigma\mapsto 1 \stackrel\sigma\mapsto 2 \stackrel\sigma\mapsto 
 3 \stackrel\sigma\mapsto 4 \stackrel\sigma\mapsto 5 \stackrel\sigma\mapsto 
 6 \stackrel\sigma\mapsto \dots  n \stackrel\sigma\mapsto \sigma(n) \dots
 \]

The first two Peano axioms (definition~\ref{def:assiomi_peano}) 
serve to ensure that in this process of \emph{counting}%
\mymargin{counting}%
\index{counting}
we always find different numbers (we never go back) since no number 
can have $0$ or a previously encountered number as a successor (which 
if not zero is the successor of another number).
This property is in some ways paradoxical
(Galileo's paradox, or Hilbert's hotel paradox)
\index{paradox!of Galileo}%
\index{paradox!of infinite sets}%
\index{Hilbert's hotel}%
\index{Hilbert!hotel}%
and for Dedekind, it is the definition of an infinite set (we will introduce it
shortly).%
\mynote{Galileo Galilei (1564--1642) see historical notes on page~\pageref{nota:Galileo}}%

The axioms of set theory that we have seen in previous chapters
do not guarantee the existence of infinite sets, and therefore if 
we want to ensure that the set of natural numbers actually exists, 
we must add an additional axiom. 

We could simply require by axiom that there exists a set $\NN$ that satisfies
Peano's axioms~\ref{def:assiomi_peano}.
For example, we could require 
that the set of Von Neumann ordinals \eqref{eq:vonNeumann} exists. 
\mynote{
If we denote by $P(\alpha)$ the predicate 
$\emptyset\in \alpha \land \forall n\colon (n\in \alpha \implies
n\cup\ENCLOSE{n}\in \alpha)$
then the axiom that guarantees the existence of the set $\omega$ of finite
ordinals
can be written in the form:
    $\exists \omega\colon (P(\omega) \land \forall \alpha\colon
  P(\alpha)\implies \omega\subset \alpha)$.
    The function $\sigma\colon\omega\to\omega$ would then be $\sigma(n) =
  n\cup\ENCLOSE n$.
}

But it seems more natural to simply require that at least one infinite set
exists,
as we will see that this is sufficient to guarantee the existence of the set of
natural numbers
(Theorem~\ref{th:esistenza_naturali}, where it is shown that every infinite set
contains a set
that satisfies Peano's axioms).

Let us therefore investigate for a moment what it means for a set to be
infinite.
To give a definition independent of the existence of natural numbers, we can use
an idea due to Dedekind:
a set $X$ is infinite if it can be put in a one-to-one correspondence 
with one of its proper subsets: 
there exists $f\colon X\to X$ such that $f$ is injective but not surjective. 
Intuitively, this cannot happen in a finite set: 
if $X$ is finite, we expect that every subset of it is strictly smaller than $X$
and therefore cannot be put in correspondence with all of $X$.

\begin{definition}[infinite set]
  \label{def:infinito}%
  We will say that a set $X$ is \emph{finite}
  \index{set!finite}%
  \index{finite!set}%
  \index{Dedekind!finite}%
  (more precisely: Dedekind-finite)
  if every injective function $f\colon X\to X$ is also surjective.
  \mymargin{finite set}%
\index{finite set}%

  We will say that a set $X$ is \emph{infinite} 
  (more precisely Dedekind-infinite)
  \index{set!infinite}%
  \index{infinite!set}%
  \index{Dedekind!infinite}%
  if it is not finite, that is,
  if there exists $f\colon X\to X$ injective but not surjective.
  \mymargin{infinite set}%
\index{infinite set}%
\end{definition}

As we have already observed, Peano's axioms (definition~\ref{def:assiomi_peano})
require that the function $\sigma\colon \NN\to\NN$ be injective but not
surjective.
It follows that $\NN$ must be an infinite set.

But it cannot be excluded that all sets are finite (and therefore $\NN$ does not
exist)
unless we introduce the following axiom.

\begin{axiom}[infinite]
  \label{axiom:infinito}%
  There exists an infinite set. 
\end{axiom}

Thanks to the previous axiom, the following theorem guarantees
the existence of a set $\NN$ that satisfies Peano's axioms.

\begin{theorem}[existence of natural numbers]
  \label{th:esistenza_naturali}%
If $X$ is any infinite set (definition~\ref{def:infinito}),
there exist $\NN\subset X$, $0\in \NN$, and $\sigma\colon \NN\to\NN$ 
that satisfy Peano's axioms (definition~\ref{def:assiomi_peano}).
\end{theorem}
%
\begin{proof}
If $X$ is infinite, there exists $f\colon X\to X$ injective but not surjective. 
Choose arbitrarily $0\in X\setminus f(X)$. 
Given a subset $I\subset X$, we will say that $I$ is \emph{inductive}
\mymargin{inductive}%
\index{inductive}%
if $0\in I \land (n\in I\implies f(n)\in I)$. 
We can then define:
\[
  \NN = \bigcap \ENCLOSE{I\subset X\colon \text{$I$ inductive}}.
\]
It is easily verified that $\NN$ is also an inductive subset of $X$.
\mynote{In definition~\ref{def:restrizione}, the symbol
$\llcorner$ is introduced for function restriction.}
Thus, for every $n\in \NN$, we have $f(n)\in \NN$, and therefore we can 
define $\sigma\colon \NN\to \NN$ as the restriction of $f$ 
to $\NN$: $\sigma=f\llcorner \NN$.

We then observe that $\sigma$ satisfies Peano's axioms.
The first axiom is a consequence of the injectivity of $f$.
The second is verified by how we chose $0$.
To verify the third axiom, consider any set 
$A\subset \NN$ such that $0\in A$ and such that if $n\in A$, then $n+1\in A$.
By definition, $A$ is inductive, and therefore certainly 
$\NN\subset A$ since $\NN$, as defined,
is a subset of every inductive set.
\end{proof}

The previous construction tells us that the set $\NN$ is the smallest infinite
set
in the sense that: if $X$ is infinite, then $X$ contains an isomorphic copy of
$\NN$.
\mynote{
Thanks to Theorem~\ref{th:cantor_bernstein}, we could say that if we have two
sets
that satisfy Peano's axioms, since there is an injective embedding of each
into the other (Theorem~\ref{th:esistenza_naturali}), 
then there exists a one-to-one correspondence between the two sets.
In Theorem~\ref{th:unicitaN}, we will see that not only does a one-to-one
correspondence exist,
but also that there exists such a correspondence that preserves the structure 
(i.e., the operation $\sigma$ and the element $0$). 
}

\subsection{principle of induction}

The last of Peano's axioms 
tells us that if a subset of the natural numbers contains zero and contains the
successor of each of its elements, then it contains all the natural numbers.
It serves to ensure that the counting process exhausts all 
natural numbers, and therefore there are no \emph{unreachable} naturals
starting from zero.
This property is usually used through the following.

\mynote{%
Formally, we should say that $P$ is a function from $\NN$ to $\{V,F\}$
where $V$,$F$ are two distinct values that we interpret as true and false.
In fact, the concept of \emph{predicate} is not defined within the 
formal system, but it is an external concept.
But it is obvious that to each predicate $P$ corresponds 
a function $P\colon \NN\to \{V,F\}$ and vice versa.
}

\index{principle!of induction}%
\index{mathematical induction}%
\begin{theorem}[principle of induction]
  Let $P(n)$ be a predicate.
  If 
  \begin{enumerate}
    \item $P(0)$ holds
    \item $\forall n\in \NN\colon P(n)\implies P(n+1)$
  \end{enumerate} 
  then $\forall n\in \NN\colon P(n)$.
\end{theorem}
%
\begin{proof}
  Consider the set $A=\{n\in \NN\colon P(n)\}$.
  Thanks to the hypotheses of the theorem, we can apply the third
  property of natural numbers to deduce that $A=\NN$.
  Thus, $P(n)$ is satisfied for every $n\in\NN$.
\end{proof}

Through the principle of induction, it is also possible 
to define functions (or operations) by induction.
Thanks to the following Theorem~\ref{th:induzione}, 
if we want to define a function $f\colon\NN\to X$,
it will suffice to declare the value of $f(0)$ and define $f(n+1)$ 
(i.e., $f(\sigma(n)))$ in terms of $f(n)$.
This method of defining a function is called 
\emph{recursive definition} (or by recurrence) as it defines 
the value of the function on the $n$-th term by recurring
to the value assigned on the previous terms.

\begin{theorem}[definition by induction]
  \label{th:induzione}%
  Let $X$ be a set, let $\alpha\in X$, and let $g\colon X\to X$ be a function.
  Then there exists a unique function $f\colon \NN \to X$ such that
  \begin{equation}\label{eq:4835628}
    \begin{cases}
      f(0) = \alpha, \\
      f(\sigma(n)) = g(f(n)).
    \end{cases}
  \end{equation}
  Thus, we will have
  \[
    f(0) = \alpha,\quad
    f(1) = g(\alpha),\quad
    f(2) = g(g(\alpha)),\quad
    f(3) = g(g(g(\alpha)))\dots
  \]
    More generally, if we have $\alpha\in X$ and a function $g\colon \NN \times
  X \to X$,
  there will exist a unique function $f\colon \NN \to X$ such that
  %
  \begin{equation}
    \begin{cases}
      f(0) = \alpha, \\
      f(\sigma(n)) = g(n, f(n)).
    \end{cases}
  \end{equation}
\end{theorem}
%
\begin{proof}
We must remember that functions $f\colon \NN \to X$ are nothing but relations 
and thus subsets of the product $\NN\times X$.
The properties~\eqref{eq:4835628} are thus written in the form 
$(0,\alpha)\in f$ and $(n,x) \in f \implies (\sigma(n),g(x))\in f$
(if $f(n)=x$ then $f(\sigma(n))=g(x)$).
The idea is therefore to take the smallest subset of $\NN\times X$ 
that can represent a function with the required properties.
Consider the family of sets:
\[
\mathcal F = \ENCLOSE{F\in \mathcal P(\NN\times X)\colon 
  (0,\alpha)\in F,\quad (n,x)\in F \Rightarrow (\sigma(n),g(x))\in F}.
\]
Clearly, $\mathcal F$ is not empty since $\NN\times X \in \mathcal F$.
We can then take its intersection and define a set $f$:
\[
  f = \bigcap_{F\in \mathcal F} F.
\]
The set $f$ we have defined represents a relation between $\NN$ and $X$.
Since $(0,\alpha)\in F$ for every $F\in \mathcal F$, it must be 
$(0,\alpha)\in f$.
Moreover, if $(n,x)\in f$, then $(n,x)\in F$ for every $F\in \mathcal F$, 
and therefore $(\sigma(n),g(x))\in F$ for every $F\in \mathcal F$,
from which $(\sigma(n),g(x))\in f$. This means that $f\in \mathcal F$.

We now want to prove that $f$ is a function, i.e., that it is uniquely defined 
on all of $\NN$.
First, consider the set on which $f$ is defined 
and that is $A=\ENCLOSE{n\in \NN\colon \exists x\in X\colon (n,x)\in f}$,
and prove, by induction, that $A=\NN$.
Indeed, $(0,\alpha)\in f$, so $0\in A$. 
And if $n\in A$, we know that there exists $x\in X$ such that $(n,x)\in f$, 
and thus, being $f\in \mathcal F$, also $(\sigma(n),g(x))\in f$,
from which $\sigma(n)\in A$. 
We have shown that $f$ is defined on all of $\NN$.

Now, let's prove that $f$ is unique. Consider 
the set on which $f$ is uniquely defined: 
$B=\ENCLOSE{n\in \NN\colon \exists! x\in X\colon (n,x)\in f}$.
Again, we want to prove by induction that $B=\NN$. 
To prove that $0\in B$, since we already know that $(0,\alpha)\in f$, 
we must prove that if $x\neq \alpha$, then $(0,x)\not\in f$.
So let $x\neq \alpha$ and consider 
the set $F=f\setminus\ENCLOSE{(0,x)}$.
Clearly, $F\in \mathcal F$ since $(0,\alpha)\in F$
because $(0,\alpha)\in f$ and $(0,\alpha)\neq (0,x)$.
Moreover, if $(n,y)\in F$, then $(n,y)\in f$, 
and therefore $(\sigma(n),g(y)) \in f$.
But certainly $(\sigma(n),g(y))\neq (0,x)$ since $\sigma(n)\neq 0$, 
thus $(\sigma(n),g(y))\in F$.
Since $F\in \mathcal F$, it must be $f\subset F$, and thus 
$(0,x)\not \in f$. Therefore, $f$ is uniquely defined at $0$.

We must now show that if $n\in B$, then also $\sigma(n)\in B$.
If $n\in B$, it means that there is a unique $x\in X$ such that $(n,x)\in f$.
Certainly, also $(\sigma(n),g(x))\in f$.
Then take $y\neq g(x)$, and we want to prove that $(\sigma(n),y)\not \in f$.
Consider, similarly to before, the set $F=f\setminus\ENCLOSE{(\sigma(n),y)}$
and try to prove that $F\in \mathcal F$.
Clearly, $(0,\alpha)\in F$ because $(0,\alpha)\in f$, 
and it cannot be $(0,\alpha)=(\sigma(n),y)$ since $\sigma(n)\neq 0$.
If we now suppose that $(m,z)\in F$, then certainly $(m,z)\in f$, 
and thus $(\sigma(m),g(z))\in f$: 
we must show that $(\sigma(m),g(z))\in F$. 
On the other hand, if it were $(\sigma(m), g(z))=(\sigma(n),y)$, 
it would have to be $m=n$ since $\sigma$ is injective. 
But since $f$ is uniquely defined on $n$, it must then also be 
$(n,z) = (n,x)$, and therefore $g(z)=g(x) \neq y$. 
Thus, $(\sigma(m),g(z))\in F$, and $F\in \mathcal F$.
But then $f\subset F$, and thus $(\sigma(n),y)\not \in f$.
This means that $f$ is uniquely defined also at $\sigma(n)$.
By induction, $B=\NN$, and $f$ is a function $f\colon \NN\to X$.

Obviously, since $f\in \mathcal F$, we know that $f$ 
satisfies the properties required by the theorem.

In the second part of the theorem, where $g\colon \NN\times X \to X$,
we can consider the set $Y=\NN\times X$ and the function 
$G\colon Y\to Y$ defined by $G(n,x) = (\sigma(n), g(n,x))$.
Then, applying the first part, we can find $F\colon \NN\to Y$
such that $F(0) = (0,\alpha)$ and $F(\sigma(n)) = G(F(n))$.
It will suffice to take as $f(n)$ the second component of $G(n)$.
\end{proof}

\begin{theorem}[addition on $\NN$]
  \label{th:addizione_naturali}%
  \index{addition!on $\NN$}%
  \mymargin{addition}%
On $\NN$, there is a uniquely defined addition operation 
that satisfies the following properties:
\begin{equation}\label{eq:3408923}
  \begin{cases}
    n + 0 = n,\\
    n + \sigma(m) = \sigma(n+m).
  \end{cases}
\end{equation}

This operation makes $\NN$ a commutative additive monoid 
with the neutral element $0$.
\end{theorem}
\begin{proof}
The addition operation is uniquely defined 
thanks to Theorem~\ref{th:induzione}.

To prove that $\NN$ is a commutative monoid, we must verify
that $0$ is the neutral element and that the associative 
and commutative properties hold.

Clearly, $n+0=n$ by definition~\eqref{eq:3408923}. 
We now verify that $0+m=m$ as well, and we do this by induction on $m$.
If $m=0$, it is true by definition.
Regarding $0+\sigma(m)$, we have by definition:
$0+\sigma(m)=\sigma(0+m)$, and by the inductive hypothesis, $0+m=m$.
Thus, $0+\sigma(m)=\sigma(m)$, which is what we wanted to prove.

We verify the associative property: $n+(m+k) = (n+m)+k$, and 
we do this by induction on $k$. 
If $k=0$, then $n+(m+0)=n+m=(n+m)+0$, by the previous properties.
For the inductive step, we evaluate (always using
definition~\eqref{eq:3408923}):
$n+(m+\sigma(k)) = n+\sigma(m+k) =\sigma(n+(m+k))$.
The inductive hypothesis guarantees that $n+(m+k)=(n+m)+k$, and therefore:
$n+(m+\sigma(k)) = \sigma((n+m)+k) = (n+m)+\sigma(k)$,
which is what we wanted to prove.

Finally, we verify the commutative property: $n+m=m+n$.
We do this by induction on $m$. 
If $m=0$, it is trivial due to the properties of the neutral element.
The inductive step is immediately 
obtained thanks to definition~\ref{eq:3408923} 
and the inductive hypothesis: 
$n+\sigma(m) = \sigma(n+m) = \sigma(m+n) = m+\sigma(n)$.
\end{proof}

\begin{example}\label{ex235}
  Prove that $2+3=5$.
\end{example}  
%
\begin{proof}[Solution]
Recall that, by definition \eqref{eq:cifre}, 
$1=\sigma(0)$, $2=\sigma(1)$, $3=\sigma(2)$, $4=\sigma(3)$, and $5=\sigma(4)$.
We need to use the properties $n+\sigma(m) = \sigma(n+m)$ 
and $n+0=n$:
\begin{align*}
2+3 &= 2 + \sigma(2) = \sigma(2+2) = \sigma(2+\sigma(1))
=\sigma(\sigma(2+1)) = \sigma(\sigma(2+\sigma(0)))\\
&=\sigma(\sigma(\sigma(2+0)))
=\sigma(\sigma(\sigma(2)))
=\sigma(\sigma(3))
=\sigma(4)=5.
\end{align*}
\end{proof}

\mymargin{$\sigma(n)=n+1$}%
Since $\sigma(n) = \sigma(n+0) = n+\sigma(0) = n+1$, 
we can avoid using the \emph{successor} function in the future 
and write directly $n+1$ to indicate the successor of $n$.

We have defined addition $n+m$ 
as the iteration of the successor function.
In general, it is possible to define multiplication by a natural number 
as a repeated sum, and exponentiation (with a natural exponent)
as a repeated multiplication.
This can be done on $\NN$ as on any monoid.

\begin{theorem}[iterated operations]
\label{th:operazioni_iterate}%
\label{th:operazione_ripetuta}%
If $M$ is an additive monoid (definition~\ref{def:monoide}),
we can define by induction the multiplication $n\cdot x$ between a natural 
number $n\in \NN$ and an element $x$ of the monoid:
\[
\begin{cases}
  0\cdot x = 0 \in M \\
  (n+1)\cdot x = n\cdot x + x.
\end{cases}
\]

Moreover, the following properties hold
(valid for every $n,m\in \NN$ and every $x,y\in M$):
\begin{enumerate}
  \item absorbing element: $0\cdot x = 0$ and $n\cdot 0 = 0$;
  \item distributive properties: 
    $n\cdot (x+y)=n\cdot x + n\cdot y$
    and $(n+m)\cdot x = n\cdot x + m\cdot x$;
  \item associative property: $n\cdot (m\cdot x)= (n\cdot m)\cdot x$.
\end{enumerate}

The same can obviously be done in multiplicative monoids
if we define the power $x^n$ as a repeated product:
\[
\begin{cases}
  x^0 &= 1 \in M \\
  x^{n+1} &= x^n \cdot x,
\end{cases}
\]
the properties are analogous:
\begin{enumerate}
  \item trivial powers: $x^0=1$, $1^n=1$;
  \item power of a product: $(x\cdot y)^n = x^n\cdot y^n$;
  \item product of powers: $x^{n+m}=x^n\cdot x^m$;
  \item power of a power: $(x^m)^n = x^{m\cdot n}$.
\end{enumerate}
\end{theorem}
\begin{proof}
The properties will be proven by induction.
Let's consider the case where $M$ is an additive monoid.
The absorbing property $0\cdot x = 0$ is given by definition,
while $n\cdot 0 = 0$ is immediately proven 
by induction; indeed, for $n=0$, we have $0\cdot 0 = 0$, 
and the inductive step is $(n+1)\cdot 0 = n\cdot 0 + 0 = 0 + 0 = 0$.

For the first distributive property: $n\cdot (x+y)=n\cdot x + n\cdot y$,
the base case is verified: $0\cdot(x+y) = 0 = 0 + 0 = 0\cdot x + 0\cdot y$.
The inductive step becomes: 
$(n+1)\cdot(x+y)
=n\cdot(x+y)+x+y
=n\cdot x + x + n\cdot y + y 
=(n+1)\cdot x + (n+1)\cdot y$.
The second $(n+m)\cdot x = n\cdot x + m\cdot x$
for $n=0$ is trivial: $(0+m)\cdot x = m\cdot x = 0\cdot x + m\cdot x$.
The inductive step:
$(n+1+m)\cdot x = (n+m)\cdot x + x 
= n\cdot x + x + m\cdot x
= (n+1)\cdot x + m\cdot x$.

The associative property: $(n\cdot m)\cdot x = n\cdot (m\cdot x)$
is obvious for $n=0$: $(0\cdot m)\cdot x = 0 = 0\cdot x = 0\cdot (m\cdot x)$.
For the inductive step, we also use the distributive property:
$((n+1)\cdot m)\cdot x 
= (n\cdot m+m)\cdot x 
= (n\cdot m)\cdot x + m\cdot x
= n\cdot (m\cdot x)+m\cdot x
= (n+1)\cdot (m\cdot x)$.

For the multiplicative monoid, the properties are exactly 
the same; only the notation with which we write them changes:
instead of the sum on $M$, we write a product, and instead of the product, 
we write a power. 
\end{proof}

\begin{exercise}
  Prove that $2\cdot 3 = 6$ and that $2^3 = 8$.
\end{exercise}
  
We define the ordering of $\NN$ 
with the idea that a number $m$ is \emph{larger} than $n$ 
if it is obtained by adding something to $n$.
\index{relation!order on $\NN$}%
\mymargin{ordering of $\NN$}%
\[
    n\le m \iff \exists k\in \NN \colon n+k = m.
\]

\begin{theorem}[properties of operations on $\NN$]
\label{th:operazioni_naturali}%
\index{multiplication!on $\NN$}%
\index{exponentiation!on $\NN$}%
The operations of addition, multiplication, exponentiation, 
and ordering defined on $\NN$
satisfy the following properties.
Properties of addition and multiplication:
\begin{enumerate}
    \item neutral element:
      $n + 0 = 0 + n = n$,
      $n\cdot 1 = 1\cdot n = n$;
    \item absorbing element:
      $n\cdot 0 = 0$;
    \item associative property: 
      $(n+m)+k = n+(m+k)$, $(n\cdot m)\cdot k = n \cdot (m\cdot k)$
    \item commutative property: 
      $n+m = m+n$, $n\cdot m = m\cdot n$;
    \item distributive property:
     $k\cdot (n+m) = k\cdot n + k\cdot m$;
    \item invariant property:
     if $m+k = n+k$, then $m=n$;
    \item cancellation of product: if $m\cdot n=0$, then $m=0$ or $n=0$;
\end{enumerate}
properties of exponentiation:
\begin{enumerate}
  \item $n^0 = 1$;
  \item $n^1 = n$;
  \item $n^{m+k} = n^m \cdot n^k$;
  \item $(n^m)^k = n^{m\cdot k}$;
  \item $(n\cdot m)^k = n^k \cdot m^k$;
\end{enumerate}
properties of ordering:
\begin{enumerate}
  \item $n\le n$ (reflexive property);
  \item if $n\le m$ and $m\le k$, then $n\le k$ (transitive property);
  \item if $n\le m$ and $m\le n$, then $n=m$ (antisymmetric property);
  \item either $n\le m$ or $m\le n$ (dichotomy);
  \item monotonicity:
   if $m\le n$, then 
   $m+k\le n+k$, $m\cdot k \le n\cdot k$,
   $k^m \le k^n$, and $m^k \le n^k$. 
\end{enumerate}
\end{theorem}
%
\begin{proof}
Many of these properties are valid in any monoid and have already 
been proven in Theorem~\ref{th:operazioni_iterate}:
neutral element, absorbing element, associative property, 
distributive property, and all properties of powers.
The fact that addition is commutative and associative has already been proven 
in Theorem~\ref{th:addizione_naturali}.
Let's prove the remaining properties.

Prove the invariant property $n+k=m+k \implies n=m$ 
by induction on $k$.
For $k=0$, the property is trivial.  
For the inductive step, observe that if $m+k+1=n+k+1$,
then, thanks to the injectivity of $\sigma$, 
it must be $m+k = n+k$,
and by the inductive hypothesis, we conclude $m=n$, as we needed to prove.
  
The property $1\cdot n=n$ is a consequence 
of the definition: $(0+1)\cdot n = 0\cdot n + n = 0 + n = n$. 
The property $n\cdot 1=n$ can be proven by induction.
If $n=0$, it is given by definition. 
For the inductive step, we have $(n+1)\cdot 1 = n\cdot 1 + 1$,
and being $n\cdot 1=n$ by the inductive hypothesis, 
we conclude $(n+1)\cdot 1 = n+1$, which is what we needed to prove.

The commutative property $n\cdot m=m\cdot n$ can be proven 
easily by induction on $n$. 
For $n=0$, it is trivial, thanks to the absorbing property.
For the inductive step, observe that $(n+1)\cdot m = n\cdot m + m$,
and supposing by the inductive hypothesis that $n\cdot m=m\cdot n$, 
and using the distributive property, we obtain:
$(n+1)\cdot m = m\cdot n + m\cdot 1 = m\cdot (n+1)$,
which is what we needed to prove.

To prove the cancellation of the product, suppose for contradiction that 
the product of two numbers different from $0$ can be zero.
A natural number different from zero is the successor of another natural number,
so we could write $(n+1)\cdot(m+1)=0$. 
But $(n+1)\cdot(m+1) = n\cdot m + n + m + 1$ is the successor of a natural
number,
so it cannot be zero.

Let's see the properties of ordering.
The fact that $n+0=n$ proves the reflexive property: $n\ge n$.
  
For the transitive property, it is sufficient to observe
that if $n=m+j$ and $m=k+l$, then $n=k+l+j$.

For the antisymmetric property, suppose we have $n=m+k$ 
and $m=n+j$. 
We deduce that $ m = m+k+j$, from which $k+j=0$.
If it were $k\neq 0$ or $j\neq 0$, the left side would be the successor 
of a natural number, but $0$ is not the successor of any natural number.
This means that $k=j=0$, and therefore $n=m$, as we wanted to prove.

To show that the ordering is total, we must 
proceed with an inductive proof. 
By induction on $n\in \NN$, we want to show 
that for every $m\in \NN$, we have $n\le m$ or $m\le n$. 
If $n=0$, the fact is obvious since $m=m+0$, and therefore $m\ge 0$
for every $m\in \NN$.
Fixed $m$ and $n$, suppose now we know that $n\le m$ or $m\le n$,
and we prove that then $n+1\le m$ or $m\le n+1$.
If $m\le n$, it means that $n = m + k$ for some 
$k\in \NN$. But then $n+1 = m + k + 1$, and therefore 
it also holds that $m\le n+1$. 
If instead $n\le m$, it means that there exists $k\in \NN$ 
such that $m=n+k$. If $k\neq 0$, then $k=1+j$ with $j\in \NN$,
from which $m=n+1+j$, and thus also $n+1\le m$.
If $k=0$, then $n=m$, and thus $n+1=m+1$, which leads us 
to the reverse inequality $m \le n+1$.
In any case, the inductive step is proven.

Now, let's prove the monotonicity of sum and product.
Let $m\le n$ and let $k$ be arbitrary.
Since $m\le n$, we have $n=m+j$. 
Thus,
\[
n+k = m + j + k \ge m+k   
\] 
and
\[
m\cdot k = (n+j)\cdot k = n\cdot k + j\cdot k \ge n\cdot k.  
\]

To verify the monotonicity of powers, we proceed by induction.
It will suffice to prove that $m^k\le(m+1)^k$ and $k^m\le k^{m+1}$.
We do this, in turn, by induction:
\[
(m+1)^{k+1} = (m+1)\cdot (m+1)^k \ge m\cdot m^k
\]
and (assuming $k\ge 1$)
\[
k^{m+1} = k\cdot k^m \ge 1\cdot k^m = k^m. 
\]
\end{proof}

\begin{exercise}[notable products]
Using the distributive property, prove that, for $a,b\in \NN$, we have:
\[
(a+b)^2 = a^2+2ab+b^2,
\qquad
(a+b)^3 = a^3 + 3a^2b + 3ab^2 + b^3.
\]
\end{exercise}

The invariant property tells us 
that if there exists $k\in \NN$ such that $n+k=m$,
then such $k$ is unique. 
\index{difference}%
\index{subtraction}%
\mymargin{difference}
Thus, if $m\ge n$, we can define the difference $k=m-n$
as that unique $k$ such that $n+k=m$.

Similarly, if there exists $k\in \NN$ such that $m=k\cdot n$,
we will say that $m$ is a \emph{multiple} of $n$, 
\index{multiple}%
\mymargin{multiples and divisors}%
or that $n$ \emph{divides} $m$ (we can write $n\vert m$).
We will say that $n$ is \emph{even} if $n$ is divisible by two,
otherwise, we will say that $n$ is \emph{odd}.
If $n\neq 0$ and $m=k\cdot n$, then $k$ is unique 
(it can be verified using the cancellation property of the product)
and is called the \emph{quotient} of $m$ divided by $n$, 
\index{quotient}%
we will write: $k=\frac m n$.

Note that we have defined $n^0=1$ for every $n\in \NN$,
including $n=0$. 
Thus, we have deliberately defined $0^0=1$:
this (controversial) definition will prove useful.
Note instead that $0^n=0$ only if $n\neq 0$.
Indeed, any number multiplied by $0$ gives zero, 
so a repeated multiplication gives zero 
if there is at least one zero factor. 
But in the product $0^0$, there are $0$ zero factors, so there is actually no 
multiplication by $0$. 
It is therefore natural that the result is $1$, 
the neutral element of multiplication.

\begin{exercise}
  Prove that for every $n\in \NN$, we have:
  \[
    n \text{ odd} \implies \exists k\in \NN\colon n=2k+1.
  \]
\end{exercise}
  
\begin{exercise}
  Prove by induction that for every $n\in \NN$, $n\ge 4$, we have 
  \[  
    2^n \ge n^2.
  \]
  In other words, prove that for every $n\in \NN$, we have 
  \[
    2^{n+4} \ge (n+4)^2.  
  \]
  \end{exercise}
  
\subsection{finite sequences or n-tuples}
\index{n-tuples}%
\index{sequence!finite}%

We have already seen that if $A$ is a set, we can define the set 
of pairs of elements of $A$ with the Cartesian product between sets: 
$A\times A$. 
What does a repeated product represent?
The sets $(A\times A)\times A$ and $A\times(A\times A)$ are not 
equal, as the first is a set of pairs whose first element 
is itself a pair, while the second set is a set of pairs 
whose second element is a pair:
\begin{gather*}
  (A\times A)\times A = \ENCLOSE{((a_0,a_1),a_2)\colon a_0,a_1,a_2\in A},
  \\
  A \times (A\times A) = \ENCLOSE{(a_0,(a_1,a_2))\colon a_0,a_1,a_2\in A}.
\end{gather*}
These sets are very similar to each other and could be identified.
Another very similar set is the set $A^{3}$.  
Using Von Neumann's definition, $3 = \ENCLOSE{0,1,2}$, the set 
$A^{3}$ represents the set of all functions 
$\vec a\colon \ENCLOSE{0,1,2}\to A$.
The function $\vec a$ is uniquely determined by its value at the 
three points of the domain: $a_0 = \vec a(0)$, $a_1=\vec a(1)$, $a_2=\vec a(2)$
since $\vec a = \ENCLOSE{0\mapsto a_0, 1\mapsto a_1, 2\mapsto a_2}$.
We can therefore identify $\vec a$ with the triple (or vector) of values:
\[
  \vec a = (a_0, a_1, a_2), \qquad a_0,a_1,a_2 \in A.  
\]
The values $a_0$, $a_1$, and $a_2$ are also called \emph{coordinates}
or \emph{components} of the \emph{vector} $\vec a$.
\mymargin{coordinates, components, vector}%
\index{coordinates, components, vector}%
\index{coordinates!vector}%
\index{components!vector}%
\index{vector}%
Indeed, the set $A^{2}$ can be identified with $A\times A$, the set 
$A^{3}$ will be the set of triples of elements of $A$, and in general, if
$n\in\NN$,
we identify $A^{n}$ with the set of $n$-tuples (read: n-tuples) of elements 
\index{n-tuples}%
of $A$:%
\mynote{
  We are here identifying $n\in \NN$ with the set $\ENCLOSE{0,1,\dots, n-1}$.}%
\[
   \vec a \in A^{n} \iff 
   \vec a = (a_0, a_1, \dots, a_{n-1}).  
\]

Historically, we have become accustomed to counting starting from $1$ instead of
$0$.
For this reason, it is usual to number the elements of an $n$-tuple with indices
ranging from $1$ to $n$ instead of from $0$ to $n-1$.
Thus, if $\vec a \in A^{3}$, it will be usual to write 
$\vec a = (a_1, a_2, a_3)$ instead of $\vec a = (a_0, a_1, a_2)$.

\subsection{sums and products with a variable number of terms}
\index{summation}%
\index{product notation}%

We now introduce the concept of iterated sum (and product) on $\NN$.
We will generally consider any monoid $A$, so 
the definitions we give will be usable on $\NN$ 
as an additive monoid (and thus we will have iterated sums)
or as a multiplicative monoid (so we will have iterated products).
Furthermore, the same definition will be usable 
on numerical sets that we have not yet defined, such as 
$\QQ$, $\RR$, and $\CC$.

In general, if $A$ is a monoid (definition~\ref{def:monoide}) additive 
(or multiplicative), we can make sums (or products) of an arbitrary 
(but finite) number of terms (or factors) of $A$.
More precisely, if $\vec a \in A^n$ 
is an $n$-tuple of elements of $A$, 
we want to define its sum and product:
\mynote{In this context, it is usual to number the elements of the 
$n$-tuple starting from index $1$ instead of index $0$.}
\begin{equation}\label{eq:6122112}
\sum_{k=1}^n a_k = a_1 + a_2 + \dots + a_n,
\qquad 
\prod_{k=1}^n a_k = a_1 \cdot a_2 \dots a_n.
\end{equation}
Formally, a definition by induction must be given. 
It will suffice to say that the sum (or product)
of zero terms (or factors) equals the neutral element 
of the operation, i.e., $0$ for the sum (and $1$ for the product).
The sum of $n+1$ terms (or $n+1$ factors) is the sum of the first $n$ 
to which we add the last term (or multiply by the last factor).
Formally:
\[
  \begin{cases}
    \displaystyle\sum_{k=1}^{0} a_k = 0, \\
    \displaystyle\sum_{k=1}^{n+1} a_k = \enclose{\sum_{k=1}^{n} a_k} + a_{n+1};
  \end{cases}  \qquad
  \begin{cases}
    \displaystyle\prod_{k=1}^{0} a_k = 1, \\
    \displaystyle\prod_{k=1}^{n+1} a_k = \enclose{\prod_{k=1}^{n} a_k} \cdot a_{n+1}.
  \end{cases}
\]
The variable $k$ that appears in formulas~\eqref{eq:6122112} is \emph{dummy}: 
any other variable that does not appear elsewhere can be used in its place.
The expressions we have written depend on $n$ but not on $k$.

It is naturally possible to also make a sum starting from an index 
different from $0$:
\[
  \sum_{k=m+1}^{m+n} a_k = \sum_{j=1}^{n} a_{m+j}.
\]
From a mnemonic point of view, we have made 
a \emph{change of variable} $k=m+j$:
for $j=1$, we find $k=m+1$, and for $j=n$, we find $k=m+n$.

Note that it results in:
\[
  x^n = \prod_{k=1}^n x
\]
since the recursive definition of power coincides 
with the recursive definition of 
the product of $n$ factors equal to $x$.

\begin{theorem}
In a commutative additive monoid, we have:
  \[
  \sum_{k=1}^n  \enclose{a_k + b_k} 
  = \sum_{k=1}^n a_k + \sum_{k=1}^n b_k.
  \]

Let $f\colon A\to A$ be an \emph{additive} function,
i.e., $f(0) = 0$ and $f(x+y)=f(x)+f(y)$.
Then 
  \[
    \sum_{k=1}^n  f(a_k) = f\enclose{\sum_{k=1}^n a_k}.
  \]

Finally, it results in 
  \[
  \sum_{k=1}^{m+n} a_k = \sum_{k=1}^m a_k + \sum_{k=1}^n a_{k+m}
  \]
or 
  \[
  \sum_{k=1}^{m+n} a_k = \sum_{k=1}^m a_k + \sum_{k=m+1}^{m+n} a_k.  
  \]
\end{theorem}
\begin{proof}
  The proofs can be carried out 
  by induction.
\end{proof}

\begin{exercise}
  \label{ex:somma_lineare}%
  Prove that 
  \[
    \sum_{k=1}^n k = \frac{n\cdot (n+1)}{2}
  \]
\end{exercise}
The result of the sum can be remembered as follows:
the average of an arithmetic progression $\frac{1+2+ \dots + n}{n}$ 
is equal to the average between the first 
and the last term of the progression: $\frac{1+n}{2}$.
\begin{proof}[Solution]
We prove it by induction.
For $n=0$, the sum is equal to $0$ by definition.
If the formula is true for a certain $n$, we have 
\[
  \sum_{k=1}^{n+1} k = \enclose{\sum_{k=1}^n k} + (n+1)
   = \frac{n\cdot(n+1)}{2} + (n+1) 
   = \frac{(n+2)(n+1)}{2}
\]
which is what we wanted to prove.
\end{proof}

\begin{exercise}
\label{ex:somma_quadrati}
  Find a formula to explicitly calculate:
  \[
    \sum_{k=1}^n k^2.
  \]
\end{exercise}
\begin{proof}[Solution.]
We want a formula similar to those proven in the previous exercises. 
This time, however, we have to find the formula ourselves.
We present a method that allows, in general, 
to find the formula to calculate $\sum_{k=1}^n k^{m+1}$
if we know the formula for $\sum_{k=1}^n k^m$.

Observe that the difference 
of two consecutive cubic terms turns out to be quadratic:
\[
(k+1)^3 - k^3 = 3 k^2 + 3k + 1.  
\]
Summing both sides of the previous equation, we obtain:
\begin{equation}\label{eq:309838}
\sum_{k=1}^n (k+1)^3 - \sum_{k=1}^n k^3 
= 3\sum_{k=1}^n k^2+3\sum_{k=1}^n k+\sum_{k=1}^n 1.
\end{equation}
On the right side, the sum we want to calculate appears 
(multiplied by 3)
along with two other sums whose values we already know. 
It will then suffice to determine the value of the left side, where 
we observe that the terms of the two sums cancel each other out, 
\mynote{%
it is called a \emph{telescopic sum}
\index{sum!telescopic}%
\index{telescopic}%
since the terms of the two summations close into each other 
like the tubes of a telescope.} %
except for the last of the first sum 
and the first of the second: 
\[
  \sum_{k=1}^n(k+1)^3 - \sum_{k=1}^n k^3 = (n+1)^3 - 1^3.
\]
Multiply equation~\eqref{eq:309838} by 2,
highlight the sum of squares, and 
use the formula $2\sum_{k=1}^n k = n(n+1)$
to obtain:
\begin{align*}
  6 \sum_{k=1}^n k^2 
  &=  2(n+1)^3 - 3 n(n+1) - 2n - 2
  = (n+1) \cdot \Enclose{2(n+1)^2 - 3n} - 2(n+1)\\
  &= (n+1)\cdot \Enclose{2n^2+4n+2 -3n - 2}
   = (n+1)\cdot \Enclose{2n^2+n} 
   = n \cdot (n+1)(2n+1)
\end{align*}
from which 
\begin{equation}\label{eq:somma_quadrati}
  \sum_{k=1}^n k^2 = \frac{n\cdot (n+1)(2n+1)}{6}.
\end{equation}
\end{proof}

\begin{exercise}
Prove that 
\[
  \sum_{k=1}^n k^3 = \frac{n^2\cdot (n+1)^2}{4}.
\]
You can use the method of the previous exercise to find an explicit formula 
or, trivially, just verify by induction that the formula is valid.
\end{exercise}

\begin{exercise}
  \label{ex:somma_geometrica}%
Note that if we take a geometric progression 
(successive powers with the same base), 
for example:
\[
  x = 3^3 + 3^4 + 3^5 + 3^6
\]
each term is equal to the previous one multiplied by $3$.
Thus, if we multiply the entire sum by the base,
\[
  3x = 3^4 + 3^5 + 3^6 + 3^7
\]
we obtain the initial sum with one more term at the end 
and one less term at the beginning. 
\mynote{This is also a \emph{telescopic sum}}
By making the difference, all terms cancel except 
the first and the last
\[
 3x - x = 3^7 - 3^3  
\]
and we can solve for $x$.

Formalize the previous reasoning to prove that 
fixed $a\neq 1$, we have for every natural $n$:
  \[
    \sum_{k=0}^n a^k = \frac{a^{n+1}-1}{a-1}.
  \] 
\end{exercise}

Denote by $\Enclose{n} = \ENCLOSE{0,1,2, \dots, n-1}$.
If $\sigma\colon \Enclose{n} \to \Enclose{n}$
is bijective (such a function is called a \emph{permutation}%
\mymargin{permutation}%
\index{permutation})
then a change of variable $k=g(j)$ can be made:
\[
    \sum_{k=0}^{n-1} a_k = \sum_{j=0}^{n-1} a_{\sigma(j)}.
\]
This equality can be easily proven in the case 
where $\sigma$ swaps only two indices, leaving all others fixed 
(transposition), and then it can be extended to all permutations
by observing that every permutation can be written as a composition 
of transpositions.

This property is essentially the commutative property of the sum 
extended to any number of terms.
Thanks to this property, we can define the sum of a function 
defined on any finite set. 
If $f\colon X \to M$
where $M$ is a monoid, 
and $\sigma\colon \Enclose{n} \to X$ is a bijection,
then we can define 
\[
  \sum_{x\in X} f(x) = \sum_{j=0}^{n-1} f(\sigma(j))  
\]
since the sum on the right side does not depend on the bijection $\sigma$ that 
we have chosen.

%%%%%%%%%%%%%%%%%%%%%%%%%%%%%%%%%

\subsection{well-ordering and uniqueness of natural numbers}

Refer to definitions~\ref{def:minimo}
for the concept of minimum, upper bound, and lower bound.
%
\begin{theorem}[principle of well-ordering]
  \label{th:buon_ordinamento}
  Let $A\subset \NN$, $A\neq \emptyset$. 
  Then $A$ has a minimum.
\end{theorem}
%
\begin{proof}
  Observe that if $n\in \NN$ is a lower bound of $A\subset \NN$, 
  then either $n\in A$ and thus $n$ is the minimum of $A$,
  or $n+1$ is also a lower bound of $A$ since 
  there are no natural numbers strictly between $n$ and $n+1$.
  \mynote{If there were a natural number between $n$ and $n+1$, 
  subtracting $n$, we would have a natural number between $0$ and $1$. 
  But there is no $x\in \NN$ such that $0<x<1$.
  Indeed, if $x\in \NN$, $x>0$, then $x\neq 0$, and therefore 
  there exists $m\in \NN$ such that $x=m+1$. 
  But then $x\ge 1$.}
  Thus, if $A$ had no minimum, by the principle of induction, 
  every $n\in\NN$ would be a lower bound of $A$.
  In such a case, $A$ would have to be empty because if there existed $a\in A$, 
  certainly $a+1$ would not be a lower bound of $A$.
\end{proof}

We now want to prove that the set of natural numbers is essentially 
unique in the sense that if there are two sets that satisfy Peano's 
axioms, then it is possible to put the elements of the two sets 
in correspondence so that zero maps to zero and corresponding numbers 
have corresponding successors.

\begin{theorem}[uniqueness of natural numbers]
  \label{th:unicitaN}%
  If $\NN$ and $\NN'$ are two sets that satisfy Peano's axioms 
  with zero $0\in \NN$ and $0'\in \NN'$ and successor functions 
  $\sigma$ on $\NN$ and $\sigma'$ on $\NN'$, then
  there exists a bijective function $f\colon \NN\to \NN'$ such that 
  $f(0) = 0'$ and $f(\sigma(n)) = \sigma'(f(n))$.
\end{theorem}
%
\begin{proof}
We can define $f$ by induction:
\[
\begin{cases}
  f(0) = 0' \\ 
  f(\sigma(n)) = \sigma'(f(n))
\end{cases}  
\]
so it remains only to prove that $f$ is a bijection.

To prove that $f$ is injective, consider the set 
\[
  A=\ENCLOSE{a\in \NN\colon \exists b\in \NN\colon b\neq a, f(a)=f(b)}.
\]
If this set is empty, then $f$ is indeed injective.
Suppose then, for contradiction, that $A$ is not empty.
In such a case, we can consider the minimum $a=\min A$ 
(thanks to Theorem~\ref{th:buon_ordinamento}).
There must then exist $b\in \NN$ such that $b\neq a$ and $f(b)=f(a)$.
Obviously, $b\in A$, and therefore $b>a$ since
$a$ is the minimum. Thus, $b>0$, and $f(b) = f(\sigma(b-1))
=\sigma'(f(b-1))$. If $a=0$, we have $f(a)=0'$, and therefore from $f(a)=f(b)$, 
we obtain that $0'$ is in the image of $\sigma'$, which is contrary 
to Peano's axioms. 
If instead $a>0$, we will have, as for $b$,
$f(a)=\sigma'(f(a-1))$, and thus $\sigma'(f(a-1)) = \sigma'(f(b-1))$.
By the injectivity of $\sigma'$, we deduce $f(a-1)=f(b-1)$, from which 
$a-1 \in A$. But this is absurd since $a$ was the minimum of $A$.

To prove that $f$ is surjective, consider the image 
$B'=f(\NN)$ and use the principle of induction on $\NN'$ 
to prove that $B'=\NN'$.
First, $0'\in B'$ since $0'=f(0)$.
If then $n'\in B'$, then there exists $n\in \NN$ such that $f(n)=n'$.
But then $f(\sigma(n))=\sigma'(f(n))=\sigma'(n')$, 
and thus also $\sigma'(n')\in B$. 
\end{proof}

\subsection{decimal representation of natural numbers}

To efficiently represent any natural number, we use the 
\emph{decimal positional notation}%
\mymargin{decimal positional notation}%
\index{notation!decimal positional}.
Recall that the decimal digits $0,1,2,\dots, 9$ were defined in~\eqref{eq:cifre}
on page~\pageref{eq:cifre}.
The number represented by a finite sequence of decimal digits can be 
defined recursively: a sequence of a single digit represents the number 
corresponding to the digit itself, a sequence of $n+1$ digits represents the 
number represented by the first $n$ digits, multiplied by $d=9+1$ (i.e., ten),
and added to the last digit. 
For example, the sequence of digits $4701$ (read: four thousand seven hundred
one)
is defined as follows:
\[ 
  4701 = ((4\cdot d+7)\cdot d+0)\cdot d+1, \qquad d=9+1.
\]
Observing that $10 = 1\cdot d + 0 = d$, it will be written:
\begin{align*}
  4701 
  & = ((4\cdot 10 + 7)\cdot 10 +0)\cdot 10 + 1 \\
    & = \mathbf 4\cdot 10^3 + \mathbf 7\cdot 10^2 + \mathbf 0\cdot 10^1 +
  \mathbf 1 \cdot 10^0.
\end{align*}

\begin{table}
  \begin{center}
    \def\tabcolsep{3pt}
    \begin{tabular}{>{\small}r|>{\small}r>{\small}r>{\small}r>{\small}r>{\small}r>{\small}r>{\small}r>{\small}r>{\small}r}
      $+$       & 1 & 2 & 3 & 4 & 5 & 6 & 7 & 8 & 9 \\ \hline
      1         & 2 & 3 & 4 & 5 & 6 & 7 & 8 & 9 & 10 \\
      2         & 3 & 4 & 5 & 6 & 7 & 8 & 9 & 10 & 11 \\
      3         & 4 & 5 & 6 & 7 & 8 & 9 & 10 & 11 & 12 \\
      4         & 5 & 6 & 7 & 8 & 9 & 10 & 11 & 12 & 13 \\
      5         & 6 & 7 & 8 & 9 & 10 & 11 & 12 & 13 & 14 \\
      6         & 7 & 8 & 9 & 10 & 11 & 12 & 13 & 14 & 15 \\
      7         & 8 & 9 & 10 & 11 & 12 & 13 & 14 & 15 & 16 \\
      8         & 9 & 10 & 11 & 12 & 13 & 14 & 15 & 16 & 17 \\
      9         & 10 & 11 & 12 & 13 & 14 & 15 & 16 & 17 & 18
      \end{tabular}
      \qquad
      \begin{tabular}{>{\small}r|r>{\small}r>{\small}r>{\small}r>{\small}r>{\small}r>{\small}r>{\small}r>{\small}r}
        $\cdot$       & 2 & 3 & 4 & 5 & 6 & 7 & 8 & 9 \\ \hline
        2         & 4 & 6 & 8 & 10 & 12 & 14 & 16 & 18 \\
        3         & 6 & 9 & 12 & 15 & 18 & 21 & 24 & 27 \\
        4         & 8 & 12 & 16 & 20 & 24 & 28 & 32 & 36 \\
        5         & 10 & 15 & 20 & 25 & 30 & 35 & 40 & 45 \\
        6         & 12 & 18 & 24 & 30 & 36 & 42 & 48 & 54 \\
        7         & 14 & 21 & 28 & 35 & 42 & 49 & 56 & 63 \\
        8         & 16 & 24 & 32 & 40 & 48 & 56 & 64 & 72 \\
        9         & 18 & 27 & 36 & 45 & 54 & 63 & 72 & 81
      \end{tabular}
    \end{center}
    \caption{The \emph{tables} of addition and multiplication.}
    \label{tab:tabelline}
\end{table}

Proceeding as in example~\ref{ex235}, one can construct the \emph{table}
of sums, where the sums of all pairs of decimal digits are reported
(table~\ref{tab:tabelline}).
Exploiting the definition and the distributive property of the product 
(Theorem~\ref{th:operazioni_naturali}),
it will not be difficult 
to construct the \emph{table} of multiplication, where the products of 
pairs of decimal digits are reported. 
For example, once it is proven that $7\cdot 7 = 49$, one can proceed 
to prove that $7\cdot 8 = 56$ as follows:
\begin{align*}
  7\cdot 8 &= 7\cdot (7+1) = 7\cdot 7 + 7 = 49+7 
  = (4\cdot 10 + 9) + 7 = 4\cdot 10 + 16 \\
  &= 4\cdot 10 + 10\cdot 1 + 6 = 5\cdot 10 + 6 = 56.
\end{align*}

Once the \emph{tables} are memorized, it is possible to perform the operations
of addition
and multiplication on any pair of natural numbers expressed in decimal form. 
The algorithm is that of column addition and multiplication that we learned in
elementary school. The correctness of these algorithms can be understood with an
example.
Let's try to prove that $34\cdot 56 = 1904$:
\begin{align*}
\mathbf{34}\cdot \mathbf{56} 
  &= \mathbf{34}\cdot (\mathbf 5\cdot 10+ \mathbf 6) 
   = \mathbf{34}\cdot \mathbf 6 + (\mathbf{34}\cdot \mathbf 5) \cdot 10 \\
    &= (\mathbf 3\cdot 10 + \mathbf 4)\cdot \mathbf 6 + (\mathbf 3\cdot 10 +
  \mathbf 4)\cdot \mathbf 5 \cdot 10 \\
    &= \mathbf{18}\cdot 10 + \mathbf{24} + (\mathbf{15}\cdot 10 +
  \mathbf{20})\cdot 10 \\
  &= \mathbf{1}\cdot 10^2+(\mathbf